% 二阶(covD)版对比,跟 w_vlad_fig.tex 是一对——同一个故事,换成外积聚合。
% 一阶故事:近重复 patch 的残差等权求和,intra-norm 修幅度修不了方向。
% 二阶故事:G = sum_p w_p x_p x_p^T 是把每个 patch 的"方向"记成一个外积再叠加。
%   扎堆的近重复 patch 会在同一个方向反复"投票",把 G 的主特征值撑得特别大,
%   画出来就是一个被拉长的椭圆,长轴指向扎堆的方向。
%   传统方法:谱指数 Sigma^0.5 把特征值压扁(椭圆变圆),但特征*向量*不变——
%            长轴还是指向那堆近重复,只是没那么夸张了(压幅度,不改方向)。
%   我们:求和之前就把每个近重复的 x_p x_p^T 按 w_p 缩小,G 本身就不会被
%        撑向重复方向,不需要再做谱压缩(alpha=1 即可)。
\begin{tikzpicture}[
  font=\small, text=ink, >=Stealth,
  ray/.style={-{Stealth[length=4.5pt]}, draw=ink!35, line width=0.55pt},
  rayD/.style={-{Stealth[length=4.5pt]}, draw=stwo!45, line width=0.55pt, densely dotted},
  ptB/.style={circle, fill=stwo!55, draw=ink!45, line width=0.4pt, minimum size=5.5pt, inner sep=0},
  ptU/.style={circle, fill=sone, draw=ink!50, line width=0.5pt, minimum size=6pt, inner sep=0},
  org/.style={cross out, draw=ink, line width=1.1pt, minimum size=5pt, inner sep=0pt},
  lbl/.style={font=\scriptsize, text=ink!75},
  tag/.style={font=\scriptsize\itshape, text=ink!60},
  panel/.style={font=\bfseries\scriptsize, text=ink!85},
]

% ====================== 左:传统 covD ======================
\begin{scope}
\node[org] (oC) at (-0.3,-0.1) {};
\node[lbl, below=2pt of oC] {origin};

\node[ptB] (b1) at (0.65,1.15) {};
\node[ptB] (b2) at (1.05,1.45) {};
\node[ptB] (b3) at (0.85,0.75) {};
\node[ptB] (b4) at (1.25,1.05) {};
\node[ptU] (u1) at (-1.55,0.65) {};
\foreach \b in {b1,b2,b3,b4} {\draw[ray] (oC) -- (\b);}
\draw[ray] (oC) -- (u1);

% 椭圆:长轴指向扎堆方向(谱压缩后仍偏,只是没那么夸张)
\draw[stwo, line width=1.1pt, rotate around={40:(oC)}]
  ($(oC)+(0,0)$) ellipse [x radius=1.55, y radius=0.55];
\node[stwo, font=\scriptsize] at (1.55,0.35) {$G^{0.5}$};

\node[tag, above=2pt of b2, xshift=6pt] {4 near-duplicates};
\node[tag, left=1pt of u1] {1 unique};
\node[panel, below=40pt of oC] {(a) traditional covD: sum, then $\Sigma^{0.5}$};
\node[lbl, below=52pt of oC, text width=5.5cm, align=center]
 {$G=\sum_p x_px_p^\top$ is stretched toward the 4 duplicates
 (they vote for the same direction 4$\times$). $\Sigma^{0.5}$ shrinks the
 \emph{eccentricity} but the major axis still points at the duplicates.};
\end{scope}

% ====================== 右:我们 ======================
\begin{scope}[xshift=7.6cm]
\node[org] (oO) at (-0.3,-0.1) {};
\node[lbl, below=2pt of oO] {origin \scriptsize(same)};

\node[ptB] (c1) at (0.65,1.15) {};
\node[ptB] (c2) at (1.05,1.45) {};
\node[ptB] (c3) at (0.85,0.75) {};
\node[ptB] (c4) at (1.25,1.05) {};
\node[ptU] (v1) at (-1.55,0.65) {};
\foreach \b in {c1,c2,c3,c4} {\draw[rayD] (oO) -- (\b);}
\draw[ray] (oO) -- (v1);

% 椭圆:更圆,不再偏向堆方向
\draw[sone, line width=1.1pt, rotate around={5:(oO)}]
  ($(oO)+(-0.1,0.05)$) ellipse [x radius=1.05, y radius=0.85];
\node[sone, font=\scriptsize] at (1.0,1.0) {$G$ ($\alpha{=}1$)};

\node[tag, above=2pt of c2, xshift=6pt] {$w{=}0.25$ each};
\node[tag, left=1pt of v1] {$w{=}1$};
\node[panel, below=40pt of oO] {(b) ours: shrink before the outer product};
\node[lbl, below=52pt of oO, text width=5.5cm, align=center]
 {$G=\sum_p w_p\,x_px_p^\top$: each duplicate's vote is cut to $1/d_p$
 \emph{before} it reinforces $G$. No spectral fix needed afterward ---
 the axis was never biased to begin with.};
\end{scope}
\end{tikzpicture}
